\section{Common crossed products and descent to factors}
\label{sec:cores}

The preceding section puts the adjusted tensor pair under one strongly
continuous normalizing unitary group. We now pass to its crossed products,
compare finite corners there, and descend the resulting conjugacy twice:
first to the adjusted tensor pair, and then through the auxiliary tensor
factor to the original Hilbert space. Only the crossed product of the algebra
whose modular group we use is initially known to be a factor.

\subsection{A concrete fixed-point description}

Let $P\subseteq\B(H)$ be a unital von Neumann algebra on a separable
Hilbert space, and let $(U_s)_{s\in\R}$ be a strongly continuous unitary
group with $U_sPU_s^*=P$.  On $\mathcal K=L^2(\R,H)$ define
\[
 (\lambda_s\xi)(t)=\xi(t-s),\qquad
 (\rho_s\xi)(t)=\xi(t+s),\qquad
 (\pi(x)\xi)(t)=U_t^*xU_t\xi(t).
\]
We use the same symbol $\pi$ for this faithful normal representation of
all of $\B(H)$.  It is a complete isometry: at every matrix level its
norm is the essential supremum of unitary conjugates of the input.
Write
\[
 C(P,U)=\bigl(\pi(P)\cup\{\lambda_s:s\in\R\}\bigr)''.
\]
This is the usual concrete crossed product for the action
$\Ad(U_s)|_P$.  For a second unital algebra normalized by the same
group, both its crossed product and the map $\pi$ are on this same
space.

\begin{lemma}[Common-action fixed points]
\label{lem:core-fixed-points}
With this notation,
\begin{equation}\label{eq:core-fixed-description}
 C(P,U)=
 \bigl(P\barotimes\B(L^2(\R))\bigr)^\alpha,
 \qquad
 \alpha_s=\Ad(U_s\otimes\rho_s).
\end{equation}
For $(Q_a\xi)(t)=e^{iat}\xi(t)$, character conjugation
$\beta_a=\Ad(Q_a)$ preserves $C(P,U)$, and
\begin{equation}\label{eq:core-character-fixed}
 C(P,U)^\beta=\pi(P).
\end{equation}
\end{lemma}

\begin{proof}
The generators $\pi(x)$ and $\lambda_r$ lie in
$P\barotimes\B(L^2(\R))$ and are fixed by every $\alpha_s$.
For the reverse inclusion, take a fixed operator $X$ in that tensor
product.  If $f,g\in C_c(\R)$, put
\[
 \lambda(f)=\int_\R f(s)\lambda_s\,ds,
 \qquad Y=\lambda(f)X\lambda(g)^*.
\]
These integrals are understood on vectors.  The operator $Y$ remains
fixed because translations commute with $U_s\otimes\rho_s$.
Define bounded operators from $\mathcal K$ to $H$ by
\[
 R_{f,t}\xi=\int_\R f(t-s)\xi(s)\,ds.
\]
They have norm at most $\norm{f}_2$ and depend continuously on $t$
in operator norm.  On compactly supported vector functions, Fubini's
theorem gives the integral kernel
\begin{equation}\label{eq:core-smoothed-kernel}
 k(t,r)=R_{f,t}XR_{g,r}^*,\qquad
 (Y\xi)(t)=\int_\R k(t,r)\xi(r)\,dr.
\end{equation}
The kernel is bounded and jointly norm continuous.  Each value belongs
to $P$: the evaluation operators intertwine the constant action of
$P'$, and $X$ commutes with that action.  Fixedness of $Y$ gives
\[
 k(t+s,r+s)=U_s^*k(t,r)U_s.
\]
For fixed $s$ the identity first holds almost everywhere by testing the
two kernels on compactly supported vectors.  Norm continuity then
makes it hold at every $(t,r)$.  Consequently
\begin{equation}\label{eq:core-kernel-covariance}
 k(t,r)=U_t^*k(0,r-t)U_t.
\end{equation}

To reconstruct a bounded operator from this kernel inside the generated
algebra, use the masks
\[
 \chi_j(v)=(1-|v|/j)_+\qquad(j\ge1).
\]
The function $\chi_j$ is the Fourier transform of the probability
measure with density
\[
 \frac{j}{2\pi}
 \left(\frac{\sin(ja/2)}{ja/2}\right)^2\,da,
\]
with the continuous value at $a=0$.  Thus multiplication of an operator
kernel by $\chi_j(t-r)$ is a contraction: it is the weak integral of
$\Ad(Q_a)$ against this measure.  Let $Y_j$ be the masked operator.
Using \eqref{eq:core-kernel-covariance} and $\rho_v\xi(t)=\xi(t+v)$,
\begin{equation}\label{eq:core-kernel-reconstruction}
 Y_j=\int_\R
       \chi_j(v)\pi(k(0,v))\rho_v\,dv\ \in C(P,U).
\end{equation}
Indeed the integrand has compact support and is strongly continuous,
since $v\mapsto k(0,v)$ is norm continuous.  Its vector integral is a
strong limit of finite sums of elements of $C(P,U)$.  Formula
\eqref{eq:core-smoothed-kernel} verifies the equality on compactly
supported vector functions.  On these tests the kernels are bounded
on a set of finite measure and $\chi_j(t-r)\to1$, so dominated
convergence gives $Y_j\to Y$ weakly.  The uniform norm bound extends
this convergence to all vector tests.  Hence $Y\in C(P,U)$.

Choose nonnegative $f_n\in C_c(\R)$ of integral one with supports
shrinking to zero.  Strong continuity of translations gives
$\lambda(f_n)\to I$ and $\lambda(f_n)^*\to I$ strongly.  Applying the
preceding conclusion with $f=g=f_n$ gives
$\lambda(f_n)X\lambda(f_n)^*\to X$ strongly.  This proves
\eqref{eq:core-fixed-description}.

For the second assertion,
\[
 \beta_a(\pi(x))=\pi(x),\qquad
 \beta_a(\lambda_s)=e^{ias}\lambda_s.
\]
Thus $\beta$ preserves $C(P,U)$.  An operator fixed by all $\beta_a$
commutes with the scalar multiplication algebra on $L^2(\R)$.
If it also belongs to $P\barotimes\B(L^2(\R))$, it is multiplication
by a bounded measurable $P$-valued field $F(t)$.  Here and below field
measurability is weak operator measurability; on bounded sets,
separability of $H$ permits countably many coefficient tests.
The remaining $\alpha$-fixedness says, for every $s$,
\[
 F(t+s)=U_s^*F(t)U_s\quad\hbox{for almost every }t.
\]
It follows that $G(t)=U_tF(t)U_t^*$ is translation invariant in
$L^\infty(\R;\B(H))$.  Such a field is almost everywhere constant.
To see this without choosing uncountably many exceptional sets, test
against a countable dense family of vector pairs.  Each scalar
coefficient is translation invariant as an $L^\infty$ class; convolving
with a smooth compactly supported approximate identity gives continuous
translation-invariant functions, hence constants.  Their local
$L^1$ convergence proves the assertion for the coefficients and then
for the field.  Since $U_t$ normalizes $P$, the constant value $x$
belongs to $P$.  Therefore $F(t)=U_t^*xU_t$ almost everywhere,
which proves \eqref{eq:core-character-fixed}.
\end{proof}

For each unitary action $V_s$ below, use the discrete invariant
mean from \Cref{lem:invariant-mean} and write
\[
 E_V(T)=m_sV_sTV_s^*.
\]
This UCP map fixes invariant operators and preserves every von
Neumann algebra invariant under the action.

\begin{lemma}[Completely bounded gap of common crossed products]
\label{lem:core-gap}
Let $P,Q\subseteq\B(H)$ be normalized by the same group $U$.
If $\dcb(P,Q)<\gamma$, then
\[
 \dcb(C(P,U),C(Q,U))\le\gamma.
\]
The conclusion holds with any strict upper bound in place of $\gamma$.
\end{lemma}

\begin{proof}
By \Cref{lem:tensor-gap}, tensoring the pair with
$\B(L^2(\R))$ preserves the upper bound.  Average on the common ambient
operator algebra using $V_s=U_s\otimes\rho_s$ in
\Cref{lem:invariant-mean}.
For a contraction $X$ in any matrix level of $C(P,U)$, choose a
contraction $Y$ in the same matrix level of
$Q\barotimes\B(L^2(\R))$ within the chosen gap bound.  The matrix
amplification of $E_V$ fixes $X$, takes $Y$ into $C(Q,U)$ by
\Cref{lem:core-fixed-points}, and preserves both the norm bound and
the approximation error.  Reverse $P,Q$.  This proves the assertion
at every matrix level with the same bound.
\end{proof}

\subsection{A finite corner and its inflation}

We shall use \Cref{lem:finite-neighbor} to check the other corner:
its proof first excludes proper isometries by approximation in a finite
algebra and then excludes nontrivial central projections by tracial
conjugacy averaging.  Thus no factor hypothesis on the other crossed
product is needed.

\begin{proposition}[Conjugacy through a common crossed product]
\label{prop:core-descent}
There are an absolute $\gamma_0>0$ and a nondecreasing function
$\zeta:(0,\gamma_0)\to(0,\infty)$, with $\zeta(\gamma)\to0$ as
$\gamma\downarrow0$, with the following property.  Suppose
$P,Q\subseteq\B(H)$ are unital von Neumann algebras on a separable
Hilbert space, $\dcb(P,Q)<\gamma<\gamma_0$, and a strongly continuous
unitary group $U$ normalizes both.  If $C(Q,U)$ is a type
$\II_\infty$ factor, then a unitary $a\in\B(H)$ satisfies
\[
 aPa^*=Q,\qquad \norm{a-I}\le\zeta(\gamma).
\]
No factoriality or modular interpretation of $C(P,U)$ is assumed.
\end{proposition}

\begin{proof}
Write $\mathcal A=C(P,U)$ and $\mathcal B=C(Q,U)$.
Use a slightly smaller gap bound if necessary so that
\Cref{lem:core-gap} gives $\dcb(\mathcal A,\mathcal B)<\gamma$.
The semifinite factor $\mathcal B$ contains matrix units
$(e_{ij})_{i,j\ge1}$ with $\sum_i e_{ii}=I$ strongly and $e=e_{11}$
finite and nonzero.  Here is the usual projection construction.
A faithful normal semifinite trace, together with separability of
the representation, gives an increasing sequence of finite-trace
projections with supremum $I$.  Split their successive differences
by trace continuity, and group the resulting pieces into countably
many projections of one fixed positive finite trace.  The infinite
total trace ensures that the grouping exhausts $I$.  Equal-trace
comparison gives the required matrix units.  Their weak closed span
$S$ is a unital copy of $\B(\ell^2)$, hence amenable.

Apply \Cref{thm:amenable} to $S\subset_\gamma\mathcal A$.
For $\gamma<1/100$ it gives a unitary $w$ with
\[
 wSw^*\subseteq\mathcal A,\qquad \norm{w-I}\le150\gamma.
\]
Put $\mathcal A_1=w^*\mathcal A w$, so
$S\subseteq\mathcal A_1\cap\mathcal B$ and
\begin{equation}\label{eq:core-aligned-gap}
 \dcb(\mathcal A_1,\mathcal B)<301\gamma=:\kappa.
\end{equation}
Compression by the common projection $e$ gives
\[
 A_0=e\mathcal A_1e|_{e\mathcal K},\qquad
 B_0=e\mathcal Be|_{e\mathcal K},\qquad
 d(A_0,B_0)<\kappa.
\]
The algebra $B_0$ is a type $\II_1$ factor.  When $\kappa<1/10$,
\Cref{lem:finite-neighbor} first makes $A_0$ finite and then
makes it a factor.  Thus \Cref{thm:finite} applies to this pair on
the separable Hilbert space $e\mathcal K$.  Choosing its orientation
appropriately gives a unitary $v_0$ on $e\mathcal K$ with
\[
 v_0A_0v_0^*=B_0,
 \qquad \norm{v_0-e}\le\eta_{\rm fin}(\kappa).
\]

Inflate this unitary using the common matrix units:
\begin{equation}\label{eq:core-inflated-unitary}
 V=\sum_{i=1}^{\infty}e_{i1}v_0e_{1i}.
\end{equation}
The sum converges strongly and is a unitary, because it is the direct
sum of copies of $v_0$ on the orthogonal spaces $e_{ii}\mathcal K$.
Moreover
\[
 \norm{V-I}=\norm{v_0-e},\qquad Ve_{ij}=e_{ij}V.
\]
For $x\in\mathcal A_1$, its transported matrix entries
$e_{1i}xe_{j1}$ belong to $A_0$.  The entries of $VxV^*$ therefore
belong to $B_0$.  Its finite matrix compressions lie in $\mathcal B$
and converge strongly to $VxV^*$; hence $V\mathcal A_1V^*\subseteq
\mathcal B$.  The reverse argument gives equality.  Thus
\begin{equation}\label{eq:core-full-unitary}
 Z=Vw^*,\qquad Z\mathcal A Z^*=\mathcal B,
 \qquad \norm{Z-I}\le150\gamma+\eta_{\rm fin}(301\gamma).
\end{equation}

It remains to descend to $H$.  Average the character action $\beta$
on the entire $\B(\mathcal K)$, obtaining a unital completely
positive map $E_\beta$ from \Cref{lem:invariant-mean}, with
$V_a=Q_a$. By \Cref{lem:core-fixed-points}, it maps
$\mathcal B$ into $\pi(Q)$ and fixes $\pi(x)$ for every
$x\in\B(H)$.  Therefore
\begin{equation}\label{eq:core-base-ucp}
 F(x)=\pi^{-1}\bigl(E_\beta(Z\pi(x)Z^*)\bigr),\qquad x\in P,
\end{equation}
defines a unital completely positive map $P\to Q$.  Since the same
map $\pi$ is completely isometric on all of $\B(H)$,
\[
 \norm{F-\id_P}_{\cb}
 \le2\norm{Z-I}
 \le300\gamma+2\eta_{\rm fin}(301\gamma).
\]
Here $\id_P$ denotes the original inclusion into $\B(H)$.
The ordinary gap $d(P,Q)$ is less than $\gamma$.  Apply
\Cref{prop:ucp}.  Its absolute thresholds hold for all sufficiently
small $\gamma$, because $\eta_{\rm fin}(t)\to0$.  Its displacement
bound supplies $\zeta(\gamma)$.  A constant multiple of
$\gamma+\eta_{\rm fin}(301\gamma)$, enlarged to a nondecreasing
function if needed, is sufficient.  This proves the proposition.
\end{proof}

\subsection{Undoing preparation and the auxiliary tensor factor}

We finish the separable factor case.  The continuous-core input used
here is the one in \Cref{sec:modular-preparation}: the crossed product
of a type $\III_1$ factor by the modular action of a faithful normal
state is a type $\II_\infty$ factor.  We apply it only to the algebra
whose modular group remains unchanged during alignment
\cite[Corollary~9.7, p.~297]{Takesaki1973}.

\begin{theorem}[Uniform stability of separably acting factors]
\label{thm:factors}
There are an absolute $t_{\rm fac}>0$ and a nondecreasing function
$\eta_{\rm fac}:(0,t_{\rm fac})\to(0,\infty)$ with
$\eta_{\rm fac}(t)\to0$ as $t\downarrow0$ such that the following
holds.  For every separable complex Hilbert space $H$ and every pair
of unital factors $M_0,N_0\subseteq\B(H)$ with the same identity and
$d(M_0,N_0)<t<t_{\rm fac}$, there is a unitary $u\in\B(H)$ with
\[
 uM_0u^*=N_0,\qquad \norm{u-I_H}\le\eta_{\rm fac}(t).
\]
The function and the threshold are independent of the factors, their
types, and their representations.
\end{theorem}

\begin{proof}
If either factor is finite, \Cref{thm:finite} applies, with the roles
reversed if necessary.  Its finite-neighbor conclusion also shows that
for a sufficiently small gap both factors are finite in that case.
We may therefore suppose that both are infinite.

Use \Cref{prop:prepared-pair}, retaining its changes of representation
explicitly.  On $H\otimes K$ it gives the fixed hyperfinite tensor
factor $D\subseteq\B(K)$ and unitaries $u_0\in\B(H)$ and
$v\in\B(H\otimes K)$ such that
\[
 M=M_0\barotimes D,\qquad
 N=v\bigl((u_0N_0u_0^*)\barotimes D\bigr)v^*,
\]
are a prepared type $\III_1$ pair, and, for absolute constants,
\[
 \norm{u_0-I}+\norm{v-I}\le C_1t,
 \qquad \dcb(M,N)<C_2t.
\]
An enlargement of $C_2$ absorbs strict upper-bound slack.
Let $U_s=\Delta_N^{is}$ be the modular group belonging to the common
vector in that preparation.  By \Cref{prop:modular-alignment}, a
unitary $q$ satisfies
\[
 \norm{q-I}\le C_3t,\qquad
 P=qMq^*\text{ is normalized by every }U_s,
 \qquad\dcb(P,N)<C_4t.
\]
The algebra $N$ and its modular group have not changed.  Consequently
$C(N,U)$ is its type $\II_\infty$ modular core.  No modular claim
about the action on $P$ is necessary.  Apply
\Cref{prop:core-descent} with $\gamma=C_4t$ to obtain
\[
 aPa^*=N,\qquad \norm{a-I}\le\zeta(C_4t).
\]

All these operations have taken place on the actual tensor product
of the original Hilbert space.  Their inverse adjustments give
\begin{equation}\label{eq:core-undo-all}
 b=(u_0^*\otimes I_K)v^*aq,
 \qquad b(M_0\barotimes D)b^*=N_0\barotimes D,
\end{equation}
and
\begin{equation}\label{eq:core-tensor-displacement}
 \norm{b-I}\le(C_1+C_3)t+\zeta(C_4t)=:r(t).
\end{equation}
The right side tends to zero universally.

Choose a unit vector $\xi\in K$, and let $\omega_\xi$ be its normal
vector state on all of $\B(K)$.  Slicing the conjugation in
\eqref{eq:core-undo-all} defines
\[
 F_0(x)=(\id\barotimes\omega_\xi)
              \bigl(b(x\otimes I_K)b^*\bigr),\qquad x\in M_0.
\]
This is a unital completely positive map into $N_0$.  The slice is
defined on the entire ambient tensor product, so comparison with
$x\otimes I_K$ gives at every matrix level
\[
 \norm{F_0-\id_{M_0}}_{\cb}\le2\norm{b-I}\le2r(t).
\]
The reverse ordinary gap is less than $t$.  A second application of
\Cref{prop:ucp} therefore gives the desired conjugating unitary on
the original $H$, with displacement at most an absolute constant
times $r(t)$, for all sufficiently small $t$.

Choose $t_{\rm fac}$ below the finite, preparation, modular-alignment,
core-descent, and two UCP thresholds just used.  Every threshold is
absolute, and the only non-linear modulus used here is
$\eta_{\rm fin}$, which is uniform and tends to zero.  The maximum of
$\eta_{\rm fin}(t)$ and a sufficiently large constant multiple of
$r(t)$, enlarged if needed to be positive and nondecreasing, defines
$\eta_{\rm fac}$.  This establishes the stated uniform quantifiers.
\end{proof}
